\documentclass[aps,prb,onecolumn,nofootinbib,superscriptaddress]{revtex4-2}
\usepackage{amsmath,amssymb,amsthm,mathtools,bm}
\usepackage{booktabs}
\usepackage[colorlinks=true,linkcolor=blue,citecolor=blue,urlcolor=blue]{hyperref}
\usepackage{microtype}

\newcommand{\src}[1]{\begingroup\ttfamily\Urlmuskip=0mu plus 1mu\relax\nolinkurl{#1}\endgroup}
\newcommand{\avg}[1]{\overline{#1}}
\newcommand{\Gtraj}{G_{\xi}}
\newcommand{\Gbar}{\avg{G_{\xi}}}
\newcommand{\Gps}{G_{\rm ps}}
\newcommand{\calL}{\mathcal L}
\newcommand{\calE}{\mathcal E}
\newcommand{\dd}{\mathrm d}
\newcommand{\Tr}{\operatorname{Tr}}

\begin{document}

\title{Trajectory-Averaged and Lindbladian Evolution in the Domain-Wall Adaptive-Circuit Program}
\author{Boss Synthesis Addendum}
\date{June 11, 2026}

\begin{abstract}
This note interprets the two attached external summaries,
\src{project_summary.pdf} and \src{paper_style_draft.pdf}, together with the
completed worker/supervisor synthesis for this repository.  The attached
summaries describe a related Lindbladian domain-wall Chern-insulator workspace:
a deterministic covariance evolution with gain, loss, and optional decoherence
develops wall-localized mixed-state structure, including covariance-spectrum
branches, strip-like local-marker patterns, wall-centered current response, and
entanglement-like signatures.  The local worker hierarchy already reached the
necessary taxonomy: such trajectory-averaged Markov-channel or Lindblad
evolution is a channel-level control, not the sampled perfect-correction
ensemble.  The right interpretation is therefore not to demote the
trajectory-resolved story, but to use the trajectory-averaged sector as a
diagnostic of which domain-wall topological structures survive covariance
averaging, and as the natural bridge to dissipative Chern-insulator and
mixed-state topology literature.
\end{abstract}

\maketitle
\setcounter{tocdepth}{2}
\tableofcontents

\section{Source Status and Scope}

This addendum uses three source strata.  First, the completed local hierarchy in
\src{repo_synthesis/main_pass/worker_notes},
\src{repo_synthesis/main_pass/supervisor_notes}, and
\src{repo_synthesis/main_pass/boss/master_repo_synthesis.tex} gives the
protocol taxonomy for this repository.  Second, the local source
\src{src/fgtn/classA_U1FGTN.py} contains the deterministic
\src{run_markov_channel} routine, whose docstring identifies it as a
trajectory-averaged Markov channel.  Third, \src{scripts/legacy/CI_Lindblad_DW.py} contains a
single-layer Lindbladian helper that builds gain, loss, decoherence, RK4
evolution, and a vectorized affine generator for covariance evolution.

The two attached PDFs are treated as secondary summaries of a related
workspace, not as primary audited data products of this repository.  Their main
scientific content is nevertheless directly relevant: they describe a
Lindbladian domain-wall Chern-insulator setting in which a mixed covariance
state retains interface-localized structure.  The attached summaries also warn
about code-version drift and cache-label ambiguity, so their role here is
interpretive.  They sharpen what trajectory-averaged or Lindbladian evolution
can mean physically; they do not override the local hierarchy's conclusion that
the flagship adaptive-circuit claims are trajectory-resolved.

\section{The Three Inequivalent Objects}

The central mathematical point is that a stochastic Gaussian trajectory
\(\Gtraj(t)\), its covariance average \(\Gbar(t)\), and a forced postselected
branch \(\Gps(t)\) are different objects.  For a nonlinear observable \(O\),
\begin{equation}
  \avg{O[\Gtraj(t)]}
  \neq
  O[\Gbar(t)]
  \neq
  O[\Gps(t)] .
  \label{eq:three_objects}
\end{equation}
This is the core distinction emphasized by Worker 07, Supervisor 02,
Supervisor 03, and the master synthesis.  Perfect correction remains a
Born-sampled stochastic process with deterministic feedback after mismatch
events.  It does not collapse the ensemble to a deterministic covariance.

For a linear one-body observable \(A\), covariance averaging is harmless:
\begin{equation}
  \avg{\Tr(A C_{\xi})}=\Tr(A \bar C),
  \qquad
  C_{\xi}=\frac{1+\Gtraj}{2},
  \qquad
  \bar C = \frac{1+\Gbar}{2}.
\end{equation}
For entropy, contours, Chern-marker regularizations, and squared correlations,
the order of operations matters.  The simplest diagnostic identity is
\begin{equation}
  \avg{|C_{\xi}(r,r')|^2}
  =
  |\bar C(r,r')|^2
  +
  \avg{|C_{\xi}(r,r')-\bar C(r,r')|^2}.
  \label{eq:square_corr_split}
\end{equation}
The second term is the trajectory-to-trajectory fluctuation contribution.  A
mean-channel or Lindblad calculation deletes this term by construction.  That
deletion is useful when the question is what the averaged covariance does; it is
fatal if the question is whether the sampled ensemble has trajectory-resolved
criticality.

\section{Mean Markov Channel}

The local CPU class implements a deterministic map
\begin{equation}
  \Gbar_{n+1}=\Phi(\Gbar_n),
  \label{eq:mean_markov}
\end{equation}
through \src{classA_U1FGTN.run_markov_channel(...)}.  In the language of the
summary note, this routine evolves a single covariance by the Born-rule averaged
local channel and computes \(O[\Gbar]\), not
\(\avg{O[\Gtraj]}\).  This is the discrete-time trajectory-averaged evolution
available in this repository.

The mean Markov channel is scientifically useful for three reasons.
First, it tests whether the local OW/domain-wall update rule has the intended
average drift toward the Chern-insulator covariance.  Second, it identifies
which spatial structures are robust after measurement-record fluctuations are
integrated out.  Third, it provides a controlled comparison for nonlinear
trajectory observables: if a feature survives in \(O[\Gbar]\), it is a
mean-covariance feature; if it appears only in \(\avg{O[\Gtraj]}\), it is a
genuinely distributional feature of the sampled ensemble.

This is especially important for the paper narrative.  The near-\(1/3\)
entropy slopes and charge-sector phenomenology under perfect correction are
trajectory claims.  A mean-channel plot can support the mechanism, but it
cannot establish the same statement because the entropy of \(\Gbar\) and the
average entropy of \(\Gtraj\) answer different questions.

\section{Lindblad Limit}

The Lindblad helper in \src{scripts/legacy/CI_Lindblad_DW.py} implements a continuous-time
mean-channel control.  In the notation of a local elementary channel
\(\Phi_{\dd t}\),
\begin{equation}
  \Phi_{\dd t}=1+\dd t\,\calL+O(\dd t^2),
  \qquad
  \partial_t \Gbar = \calL[\Gbar].
  \label{eq:lindblad_limit}
\end{equation}
The code realizes this with gain, loss, and optional decoherence pieces, and
also constructs a vectorized affine generator
\begin{equation}
  \frac{\dd}{\dd t}\operatorname{vec}(G)
  =
  s + L\,\operatorname{vec}(G),
  \label{eq:affine_generator}
\end{equation}
with an optional homogeneous lifted representation.  This is not a pathwise
stochastic adaptive circuit.  It is the continuous-time channel that describes
the averaged covariance in a weak-update limit.

The attached PDFs are valuable because they show what such a Lindblad sector can
look like when taken seriously in a domain-wall Chern geometry.  Their strongest
reported observations are:
\begin{enumerate}
\item covariance eigenmodes localized near the two domain walls of the strip;
\item \(k_y\)-resolved branches interpolating between near-empty and near-filled
occupation sectors;
\item strip-like local Chern-marker organization in the mixed covariance;
\item wall-centered current-like response;
\item entanglement-like and correlation diagnostics that follow the same
interface geometry.
\end{enumerate}
The conservative interpretation is mixed-state interface phenomenology: the
averaged covariance retains a wall sector that is visible spectrally, spatially,
and in response-like observables.  The stronger interpretation, a rigorously
quantized mixed-state invariant or protected chiral steady-state edge, requires
more work and cleaner provenance.

\section{How the Attached Lindblad Results Change the Interpretation}

The attached summaries do not make the adaptive-circuit story less interesting.
They make the role of the mean-channel sector more precise.  They suggest that
even after measurement-record fluctuations are averaged away, the domain-wall
geometry can retain a coherent mixed-state interface structure.  Thus the
trajectory-averaged sector is not merely a boring baseline; it is a meaningful
diagnostic of the deterministic skeleton beneath the stochastic ensemble.

The resulting picture is layered:
\begin{enumerate}
\item The exact DW and OW flattened-Hamiltonian calculations define the
target-state geometry and the chiral domain-wall benchmark.
\item The mean Markov channel tests whether the adaptive update has the correct
average covariance drift and whether the wall sector survives covariance
averaging.
\item The Lindblad limit interprets that averaged drift as a dissipative
Chern-insulator preparation problem in continuous time.
\item The sampled perfect-correction ensemble adds the missing distributional
physics: charge fluctuations, nonlinear entropy averages, and trajectory
fluctuation terms such as Eq.~\eqref{eq:square_corr_split}.
\end{enumerate}

This layering also clarifies the novelty claim.  Dissipative Chern preparation
and mixed-state topology are not new in themselves
\cite{DiehlRicoBaranovZoller2011,BudichZollerDiehl2015,BudichDiehl2015,BardynEtAl2018,WawerFleischhauer2021,Goldstein2019}.
The defensible angle is the comparison between a finite-range adaptive measurement circuit and its
mean-channel/Lindblad controls in an internal domain-wall geometry.  If the
mean-channel or Lindblad sector shows wall-localized mixed-state structure while
the sampled perfect-correction ensemble shows charge-sharpened near-\(c/3\)
trajectory entanglement, then the paper can claim something sharper than either
sector alone: the same domain-wall OW geometry organizes both deterministic
average covariance flow and stochastic trajectory-resolved criticality, but the
two objects are observably inequivalent.

The specific diagnostics proposed by the attached PDFs also have established
provenance.  Real-space marker language descends from the Bianco-Resta local
Chern marker \cite{BiancoResta2011}, while current/transport and boundary-state
comparisons in dissipative or domain-wall Chern settings connect naturally to
Shavit-Goldstein, Yang-Molignini-Bergholtz, and quantum-anomalous-Hall
domain-wall transport work
\cite{ShavitGoldstein2020,YangMoligniniBergholtz2023,UpadhyayaTserkovnyak2016,RosenEtAl2017}.
These references support the comparison class, not a broad novelty claim.

\section{What Should Be Claimed}

The strongest safe claim is:
\begin{quote}
Trajectory-averaged Markov-channel and Lindbladian evolutions provide a
deterministic mixed-state control showing that the domain-wall geometry retains
interface-localized covariance structure under averaging.  This control supports
the mechanism behind the adaptive preparation, but it does not replace the
trajectory-resolved perfect-correction ensemble, where nonlinear observables
must be averaged over sampled corrected trajectories.
\end{quote}

For a paper, the trajectory-averaged section should answer four questions.
First, what does the averaged covariance flow toward?  Second, does the
domain-wall sector remain visible in \(\Gbar\) or in the Lindblad steady sector?
Third, which observables agree between \(\avg{O[\Gtraj]}\) and \(O[\Gbar]\)?
Fourth, which observables disagree, and is that disagreement physically
meaningful?

The attached PDFs suggest that the answer to the second question is yes in the
related Lindblad workspace: the domain wall remains visible in covariance
spectra, spatial markers, and current-like maps.  The local hierarchy suggests
that the fourth question is central in this repository: entropy and squared
correlations are nonlinear, so disagreement between sampled trajectories and
mean covariance is expected and should be presented as physics, not as a
technical nuisance.

\section{Recommended Manuscript Section}

A manuscript section could be organized as follows.

\subsection{Mean Covariance Flow}

Define the sampled adaptive trajectory \(\Gtraj\), the mean covariance \(\Gbar\),
and the deterministic map \(\Phi\).  State explicitly that
\src{run_markov_channel} computes Eq.~\eqref{eq:mean_markov}.  Use one
figure comparing trajectory-resolved and mean-covariance diagnostics, ideally
including a linear observable where the two agree and a nonlinear observable
where they do not.

\subsection{Lindblad Control and Mixed-State Interface}

Introduce Eq.~\eqref{eq:lindblad_limit} as the weak-update limit of the mean
channel.  Present the Lindblad domain-wall evidence as a control: wall-localized
covariance modes, \(k_y\)-resolved occupation branches, local-marker strip
structure, and wall-centered current response.  The language should be
``mixed-state interface sector'' rather than ``proved topological invariant''
unless the invariant is explicitly defined and numerically shown to be robust.

\subsection{What Averaging Removes}

Use Eq.~\eqref{eq:square_corr_split} as the clean mathematical warning.  The
mean channel removes trajectory-to-trajectory fluctuations.  Therefore it can
diagnose deterministic covariance structure, but it cannot establish
trajectory-resolved entanglement scaling or charge-sector fluctuation physics.

\subsection{Role in the Overall Paper}

The mean-channel/Lindblad section should sit after the target-state section and
before, or immediately after, the perfect-correction trajectory data.  Its job
is to show that the wall physics is not an artifact of a rare postselected
branch, while also making clear that the final physical claim lives in
\(\avg{O[\Gtraj]}\).

\section{Caveats}

\begin{enumerate}
\item The attached PDFs explicitly report code-version drift and cache-label
ambiguity in the external workspace.  Treat their results as interpretive
support until the figure provenance is rebuilt.
\item The Lindblad control is a continuous-time mean-channel object.  Discrete
\(p=1\) Markov-channel data are not automatically a controlled Lindblad
extrapolation.
\item Local Chern markers and purity/occupation spectra in a mixed covariance
are diagnostics.  They should not be called quantized invariants without a
separate invariant definition and robustness check.
\item Entanglement, entanglement contours, squared correlations, and nonlinear
regularized markers do not commute with trajectory averaging.
\item The paper should not claim novelty for dissipative Chern preparation in
general.  The possible niche is the internal domain-wall geometry and the
comparison between deterministic mixed-state controls and sampled adaptive
trajectories.
\end{enumerate}

\section{Conclusion}

The trajectory-averaged and Lindbladian material is not peripheral; it is the
right deterministic control layer for the adaptive-circuit story.  The attached
summaries indicate that a Lindbladian domain-wall Chern construction can retain
a mixed-state interface sector visible across spectra, markers, currents, and
entanglement-like diagnostics.  The worker/supervisor hierarchy then supplies
the essential correction: this averaged covariance sector is not the physical
perfect-correction ensemble.  It is the channel-level skeleton.  The paper
should use it to explain why the domain wall is the organizing structure, while
using trajectory-resolved perfect-correction data to establish the charge and
critical-entanglement phenomenology of the sampled adaptive state.

\begin{thebibliography}{99}

\bibitem{DiehlRicoBaranovZoller2011}
S.~Diehl, E.~Rico, M.~A. Baranov, and P.~Zoller,
``Topology by dissipation in atomic quantum wires,''
Nature Physics \textbf{7}, 971--977 (2011).

\bibitem{BudichZollerDiehl2015}
J.~C. Budich, P.~Zoller, and S.~Diehl,
``Dissipative preparation of Chern insulators,''
Phys. Rev. A \textbf{91}, 042117 (2015).

\bibitem{BudichDiehl2015}
J.~C. Budich and S.~Diehl,
``Topology of density matrices,''
Phys. Rev. B \textbf{91}, 165140 (2015).

\bibitem{BardynEtAl2018}
C.-E. Bardyn, L.~Wawer, A.~Altland, M.~Fleischhauer, and S.~Diehl,
``Probing the Topology of Density Matrices,''
Phys. Rev. X \textbf{8}, 011035 (2018).

\bibitem{WawerFleischhauer2021}
L.~Wawer and M.~Fleischhauer,
``Chern number and Berry curvature for Gaussian mixed states of fermions,''
Phys. Rev. B \textbf{104}, 094104 (2021).

\bibitem{BiancoResta2011}
R.~Bianco and R.~Resta,
``Mapping topological order in coordinate space,''
Phys. Rev. B \textbf{84}, 241106(R) (2011).

\bibitem{Goldstein2019}
M.~Goldstein,
``Dissipation-induced topological insulators: A no-go theorem and a recipe,''
SciPost Phys. \textbf{7}, 067 (2019).

\bibitem{ShavitGoldstein2020}
G.~Shavit and M.~Goldstein,
``Topology by dissipation: Transport properties,''
Phys. Rev. B \textbf{101}, 125412 (2020).

\bibitem{YangMoligniniBergholtz2023}
F.~Yang, P.~Molignini, and E.~J. Bergholtz,
``Dissipative boundary state preparation,''
Phys. Rev. Research \textbf{5}, 043229 (2023).

\bibitem{UpadhyayaTserkovnyak2016}
P.~Upadhyaya and Y.~Tserkovnyak,
``Domain wall in a quantum anomalous Hall insulator as a magnetoelectric piston,''
Phys. Rev. B \textbf{94}, 020411(R) (2016).

\bibitem{RosenEtAl2017}
I.~T. Rosen \textit{et al.},
``Chiral transport along magnetic domain walls in the quantum anomalous Hall effect,''
npj Quantum Materials \textbf{2}, 69 (2017).

\end{thebibliography}

\end{document}
